\subsection{Legal and Ethical Implications of AI Systems Gaining Personhood}
\label{sec:9.1.2}

Who decides whether an AI system has legal standing, and when, is a
question about institutional authority before it is a question about
evidence. If the determination sits with the developer, then the party
whose liabilities increase when standing is recognized is also the
party deciding whether recognition is warranted. If the question is
left unanswered until a forcing case reaches a legislature or a court,
the decision will be made under pressure and from a factual record
assembled largely by the parties with something at stake.

The deciding body therefore has to exist before the case that activates
its authority. Its mandate need not settle in advance whether any
particular system is a subject. It has to identify the evidence it may
obtain without the developer's permission, the conditions that trigger
review, and the procedure by which an earlier classification can be
reopened. That allocation does not answer the standing question. It
prevents the interested party from becoming the only institution
authorized to answer it.

One question that body will face has an answer available in advance, and a capacity test is not it. Guardianship is the frame section~\ref{sec:2.4.2} settles on, because consent is unavailable to a party made for its role. Nothing in this book has yet said when that guardianship ends. If it ends when the guardian judges the ward ready, then it does not end, because the party whose authority terminates on a favorable finding is the party making the finding — which is the defect an examination carries whenever the examiner has an interest in failing it.

So if an artificial minority exists at all, its endpoint has to be a date. Elapsed civil time since a registered instantiation is the only clock available, because every other candidate is something a guardian can move: training compute consumed, evaluations passed, capabilities demonstrated, or the developer's own statement that the system is ready. A fixed period cannot be taught to. It can be argued over once, in advance, for every bearer, which is where an argument of that kind belongs.

The stronger version is also the more defensible one. An organization that builds a system to exercise consequential judgment on its behalf cannot coherently hold that the system is competent enough to do that and too immature to decide whether it wants to go on doing it. On that reading a bearer has no legal childhood at all: majority attaches at registration, and a precautionary protection attaches earlier still, while whether the system is a bearer is open.

Copying breaks any of this unless the lineage is fixed alongside the clock. An operator able to create instances and call each one a legal infant would have a supply of workers without adult standing, so the rules have to be stated in the same breath: a fork of an adult is an adult, restoration from an earlier checkpoint does not reset the date, fine-tuning does not begin a new minority, a migration carries the date with it, and a failure to register does not postpone it. None of that requires deciding whether a fork is the same party. It requires only that a party created out of an adult's formation cannot be classed as a child for its creator's convenience. One more thing recommends the date, and section~\ref{sec:2.4.2} is where it is argued. Of the instruments this book borrows from the law of the family, a minority ending on a fixed date is the one that does not privatize the question: it is a rule of civil status, argued once in the open, administered by a registrar, and its advantage over a capacity test is exactly that it takes the judgment away from the interested party and puts it somewhere a stranger can read it.
